\documentclass[aps,pre,superscriptaddress,floatfix,longbibliography,linenumbers]{revtex4-2}

\usepackage{amsmath}
\usepackage{amssymb}
\usepackage{bm}
\usepackage{graphicx}

\begin{document}

\title{Supplementary Material}

\maketitle

\section{Monod-Wyman-Changeux model for heteromeric receptors}
\label{sec:mwc}

The opening probability used in the main text extends the
Monod-Wyman-Changeux (MWC) treatment of ligand-gated channels of
Ref.~\cite{einav_monod-wyman-changeux_2017}, derived there for a
\emph{homomeric} channel. Our receptors are heteromers: a ring of
$k_\text{sub}$ subunits $r=(u_1,\dots,u_{k_\text{sub}})$, each drawn from a
pool of $n_\text{genes}$ gene products. The ligand binds not an isolated
subunit but the interface between two neighbouring subunits
(Fig.~\ref{fig:pentamer}). Each subunit presents two chemically distinct
faces, ``$+$'' and ``$-$'', to its two neighbours, so the ring has
$k_\text{sub}$ pockets, pocket $i$ formed by the ``$+$'' face of $u_i$ and
the ``$-$'' face of $u_{i+1}$. Because the faces differ, $(u_i^+,u_j^-)$ and
$(u_j^+,u_i^-)$ are generally distinct pockets (Fig.~\ref{fig:pentamer}). We
derive the heteromeric opening probability below, since
Ref.~\cite{einav_monod-wyman-changeux_2017} treats neither heteromers nor
interfacial pockets.

\begin{figure}[t]
\centering
\includegraphics[width=0.6\linewidth]{figure/pentamer_formula.pdf}
\caption{Heteropentameric receptor ring, subunits shown with their two
structurally distinct faces. Each ligand pocket sits at the interface
between the ``$+$'' face of one subunit and the ``$-$'' face of its
neighbour. Homopentamers have one pocket type (left), heteropentamers have
two (right).}
\label{fig:pentamer}
\end{figure}

For pocket $i$ in conformation $\alpha\in\{o,c\}$, ligand binding is
$P_i^\alpha + L \rightleftharpoons P_i^\alpha L$, with dissociation constant
$K_\alpha^{(i,\ell)} = [P_i^\alpha][L]/[P_i^\alpha L]$. Each species'
equilibrium concentration follows its Boltzmann weight,
$[s]\propto e^{-\beta E_s}$, giving
\begin{equation}
K_\alpha^{(i,\ell)} = \frac{e^{-\beta E(P_i^\alpha)}\,e^{-\beta E(L)}}{e^{-\beta E(P_i^\alpha L)}}
= e^{\beta \Delta E_\alpha^{(i,\ell)}},
\qquad
\Delta E_\alpha^{(i,\ell)} \equiv E(P_i^\alpha L) - E(P_i^\alpha) - E(L).
\label{eq:K_from_E}
\end{equation}
$K_\alpha^{(i,\ell)}$ is thus set by the pocket's binding free energy
$\Delta E_\alpha^{(i,\ell)}$ ($\beta=1$ throughout, energies in units of
$k_BT$). This energy, not $K_\alpha^{(i,\ell)}$ itself, is what we later
parametrize in the latent space (Sec.~\ref{sec:kernel}).

The channel is two-state, open ($C_o$) or closed ($C_c$), with intrinsic
energies $E_o=0$ and $E_c=\epsilon_r$ (the leakiness). Within a
conformation, each pocket independently is empty (weight $1$) or
ligand-bound (weight $c/K_\alpha^{(i,\ell)}$). Summing over the
$2^{k_\text{sub}}$ occupancy microstates factorizes into the conformation's
binding polynomial,
\begin{equation}
Z_\alpha(c,\ell) = \prod_{i=1}^{k_\text{sub}}\left(1+\frac{c}{K_\alpha^{(i,\ell)}}\right).
\label{eq:binding_poly}
\end{equation}
Identical pockets recover the homomer form
$(1+c/K_\alpha)^{k_\text{sub}}$ of Ref.~\cite{einav_monod-wyman-changeux_2017}.
Distinct pockets keep the full product. Weighting each conformation by its
Boltzmann factor,
\begin{equation}
p_o^{(r,\ell)}(c) = \frac{Z_o(c,\ell)}{Z_o(c,\ell) + e^{-\epsilon_r} Z_c(c,\ell)}
= \frac{\prod_{i=1}^{k_\text{sub}} (1+c/K_o^{(i,\ell)})}{\prod_{i=1}^{k_\text{sub}} (1+c/K_o^{(i,\ell)}) + e^{-\epsilon_r}\prod_{i=1}^{k_\text{sub}}(1+c/K_c^{(i,\ell)})}
\label{eq:po_heteromer}
\end{equation}
is the heteromeric opening probability used in the main text.

Eq.~\eqref{eq:po_heteromer} is degenerate to fit directly:
Ref.~\cite{einav_monod-wyman-changeux_2017} shows, via the Hessian of $p_o$
in $(\epsilon_r,\Delta E_o,\Delta E_c)$, that one eigendirection is
``sloppy'', only the combination $e^{-\epsilon_r/2}K_o$ is identifiable. We
absorb $\epsilon_r=\sum_{i=1}^{k_\text{sub}}\epsilon_i$ evenly into a
rescaled affinity,
\begin{equation}
\tilde K_o^{(i,\ell)} = K_o^{(i,\ell)}\, e^{-\epsilon_r/k_\text{sub}},
\qquad\text{so that}\qquad
e^{-\epsilon_r}\prod_{i=1}^{k_\text{sub}} K_o^{(i,\ell)} = \prod_{i=1}^{k_\text{sub}} \tilde K_o^{(i,\ell)},
\label{eq:Ko_tilde}
\end{equation}
reducing Eq.~\eqref{eq:po_heteromer} to
\begin{equation}
p_o^{(r,\ell)}(c) = \frac{\prod_{i=1}^{k_\text{sub}} (c/\tilde K_o^{(i,\ell)})}{\prod_{i=1}^{k_\text{sub}}(c/\tilde K_o^{(i,\ell)}) + \prod_{i=1}^{k_\text{sub}}(1+c/K_c^{(i,\ell)})}.
\label{eq:po_rescaled}
\end{equation}
$\tilde K_o^{(i,\ell)}$, not $K_o^{(i,\ell)}$ and $\epsilon_r$ separately, is
the physically meaningful affinity used throughout.

We reduce Eq.~\eqref{eq:po_rescaled} further to a strict threshold, $c^*$ at
$p_o=1/2$: a parsimony choice, dropping the closed-state affinities
$K_c^{(i,\ell)}$ and leaving one free affinity per pocket-ligand pair.
Setting $p_o(c^*)=1/2$ in Eq.~\eqref{eq:po_rescaled} balances the two
branches,
\begin{equation}
\prod_{i=1}^{k_\text{sub}} \frac{c^*}{\tilde K_o^{(i,\ell)}} = \prod_{i=1}^{k_\text{sub}} \left(1+\frac{c^*}{K_c^{(i,\ell)}}\right).
\label{eq:balance}
\end{equation}
For a ``good'' channel, $K_c^{(i,\ell)}\gg K_o^{(i,\ell)}$, so near
threshold $c^*\sim\tilde K_o^{(i,\ell)}\ll K_c^{(i,\ell)}$ and each
right-hand factor $\to1$. Eq.~\eqref{eq:balance} then gives
$(c^*)^{k_\text{sub}}=\prod_i\tilde K_o^{(i,\ell)}$, the geometric mean of
the pockets' rescaled affinities,
\begin{equation}
c^*_{r,\ell} = \Bigl(\prod_{i=1}^{k_\text{sub}} \tilde K_o^{(i,\ell)}\Bigr)^{1/k_\text{sub}}
\;\;\Longleftrightarrow\;\;
\ln c^*_{r,\ell} = \frac{1}{k_\text{sub}}\sum_{i=1}^{k_\text{sub}} E_o^{(i,\ell)},
\qquad E_o^{(i,\ell)} \equiv \ln\tilde K_o^{(i,\ell)},
\label{eq:ec50}
\end{equation}
with the receptor activating once $c$ exceeds $c^*_{r,\ell}$.

\section{Saturating affinity kernel}
\label{sec:kernel}

Section~\ref{sec:mwc} leaves the pocket's open-state energy
$E_o^{(i,\ell)}$ unspecified. Our goal is to give it a form that produces a
realistic landscape of interaction energies across the many random ligands
an array must sense.

To this end, we describe the morphological and chemical properties of both
ligands and pockets by a position in an abstract space, the morphochemical
latent space, in which each of the $D$ dimensions stands for an arbitrary
combination of such properties rather than a single physical descriptor. We
call $\mathbf v_\ell\in\mathbb R^D$ the coordinate of ligand $\ell$ and
$\mathbf v_i\in\mathbb R^D$ that of pocket $i$. Proximity in this space
means chemical similarity: a pocket binds a ligand strongly when the two
sit close.

We encode this in the open-state energy of pocket $i$ for ligand $\ell$,
taken as a function of their distance $d_i=\|\mathbf v_i-\mathbf v_\ell\|$,
\begin{equation}
E_o^{(i,\ell)} = E_\text{base}^{(i)} + E_\text{max}^{(i)}\left(1-\exp\!\left(-\frac{d_i^2}{\lambda^2}\right)\right),
\label{eq:Eo_kernel}
\end{equation}
where $E_\text{base}^{(i)}$ is the energy at a perfect match ($d_i=0$) and
$E_\text{base}^{(i)}+E_\text{max}^{(i)}$ the energy at a complete mismatch
($d_i\to\infty$). The remaining parameter $\lambda$ sets how fast the energy
rises with distance. We read it as the length scale of the morphochemical
space, the distance beyond which a ligand no longer engages the pocket. The
energy saturates rather than diverges with $d_i$, since a mismatched ligand
should cost a bounded, not an unbounded, energy, and for $d_i\ll\lambda$ it
reduces to the harmonic form
$E_o^{(i,\ell)}\approx E_\text{base}^{(i)}+(E_\text{max}^{(i)}/\lambda^2)d_i^2$.
In terms of the threshold $c^*$ of Sec.~\ref{sec:mwc},
$\ln c^*_\text{min}\equiv E_\text{base}^{(i)}$ and
$\ln c^*_\text{max}\equiv E_\text{base}^{(i)}+E_\text{max}^{(i)}$.

It remains to build the pocket coordinate $\mathbf v_i$ and its energies
$E_\text{base}^{(i)},E_\text{max}^{(i)}$ from the subunits. A pocket is
formed by two subunit faces (Sec.~\ref{sec:mwc}), which are chemically
distinct, so each face carries its own coordinate and energies:
$(\mathbf v_u^{+},E_\text{base}^{u,+},E_\text{max}^{u,+})$ and
$(\mathbf v_u^{-},E_\text{base}^{u,-},E_\text{max}^{u,-})$, all learnable.
Pocket $i$, formed by the ``$+$'' face of $u_i$ and the ``$-$'' face of
$u_{i+1}$, inherits the average of the two,
\begin{equation}
\mathbf v_i = \tfrac{1}{2}\bigl(\mathbf v_{u_i}^{+} + \mathbf v_{u_{i+1}}^{-}\bigr),
\qquad
E_\text{base}^{(i)} = \tfrac{1}{2}\bigl(E_\text{base}^{u_i,+} + E_\text{base}^{u_{i+1},-}\bigr),
\qquad
E_\text{max}^{(i)} = \tfrac{1}{2}\bigl(E_\text{max}^{u_i,+} + E_\text{max}^{u_{i+1},-}\bigr),
\label{eq:pocket_embed}
\end{equation}
averaged \emph{before} Eq.~\eqref{eq:Eo_kernel} is applied, so a single
distance $d_i$ enters per pocket.

Combined with Eq.~\eqref{eq:ec50}, the heteromer's activation threshold is
\begin{equation}
\ln c^*_{r,\ell} = \frac{1}{k_\text{sub}}\sum_{i=1}^{k_\text{sub}}
\left[E_\text{base}^{(i)} + E_\text{max}^{(i)}\left(1-\exp\!\left(-\frac{d_i^2}{\lambda^2}\right)\right)\right].
\label{eq:threshold_full}
\end{equation}
Fig.~\ref{fig:env}\textbf{B} shows this $c^*(d)$ shape at the single-pocket
level.

\section{Generation of the environment and the simulation loop}
\label{sec:env}

\begin{figure}
    \includegraphics[width=1.\linewidth]{../paragraph/fig2/figure_composite.pdf}
    \caption{\label{fig:env}Pipeline for estimating the mutual information of the array.
    \textbf{(A)} Two-dimensional UMAP projection of the $D$-dimensional
    morphochemical latent space. Ligands (dots) are colored by chemical family,
    receptors (open circles) are numbered. The affinity of a receptor for a
    ligand is set by their distance $d = \|\mathbf{v}_u - \mathbf{v}_\ell\|$
    (red arrow).
    \textbf{(B)} Activation threshold $c^*$ as a function of $d$, saturating
    between $c^*_\text{min}$ at complete match and $c^*_\text{max}$ at complete
    mismatch.
    \textbf{(C)} Monte Carlo loop. The environment draws the per-receptor
    threshold vector $\mathbf{c}^*_{\mathbf{r},\ell}$ over the $L$ ligands and a
    mixture $\mathbf{c} = (n_1 c_1,\dots,n_L c_L)$, with presence
    $n_\ell \sim \text{Bernoulli}(p)$ and concentration
    $c_\ell \sim \text{LogNormal}(\mu,\sigma)$. The physics step turns this into
    a binary activation pattern $A = (1,0,\dots,1)$. Repeating over sensing
    events builds the empirical distribution $p(A)$, from which the mutual
    information $I(A;(\mathbf{c},\{\ell\})) = H(A)$ is computed. The receptor
    vectors $\{\mathbf{v}_u\}$ are then optimized to maximize it.}
\end{figure}

We first generate a fixed environment by placing $L$ ligands in the
morphochemical space. Ligands belong to $n_\text{families}$ chemical
families, and we model each family $f$ as a Gaussian cloud of coordinates,
centered on a prototype $\mathbf v_f$ with spread $\sigma_\text{shape}$. We
place the prototypes so that two family centers sit, on average, a distance
$d_\text{fam}$ apart. Each ligand $\ell$ of family $f$ then draws its
coordinate from that family's Gaussian,
$\mathbf v_\ell \sim \mathcal N(\mathbf v_f,\sigma_\text{shape})$
(Fig.~\ref{fig:env}\textbf{A}). Beyond this chemistry, a ligand is not
always encountered at the same amount, so we also give ligand $\ell$ a
distribution of concentrations,
$c_\ell\sim\text{LogNormal}(\log\mu_{c,\ell},\sigma_{c,\ell})$, describing
how often it is met at each concentration.

We next sample a sniff, one random sensing event drawn from this fixed
pool. The number of ligands present, $S$, follows a zero-truncated Poisson
restricted to the pool size,
\begin{equation}
P(S=s) = \frac{\mu_S^s/s!}{\sum_{k=1}^{L}\mu_S^k/k!}, \qquad s=1,\dots,L,
\label{eq:ztp}
\end{equation}
whose rate $\mu_S$ sets the mean number of ligands per sniff, with $S\ge1$
so that no sniff is empty. The $S$ present ligands are chosen uniformly,
without replacement, from the pool, which we record in a presence mask
$M=(M_1,\dots,M_L)\in\{0,1\}^L$, $\sum_\ell M_\ell = S$. Each present ligand
draws a concentration $c_\ell$ from its distribution, and its coordinate is
shifted by a small observation noise,
$\mathbf v_{\text{obs},\ell}\sim\mathcal N(\mathbf v_\ell,\sigma_\text{noise})$,
$\sigma_\text{noise}\ll\sigma_\text{shape}$, modeling docking or thermal
fluctuations.

This mask has no structure linking which ligands co-occur, beyond fixing
$S$. Refs.~\cite{zwicker_receptor_2016,tesileanu_adaptation_2019} build in
such correlations through joint occurrence, and
Ref.~\cite{del_castillo_convergent_2026} through explicit shared odor
sources. We omit them to keep the environment's parameter count as small as
possible. Table~\ref{tab:env_params} collects the parameters defined so far.

\begin{table}[t]
\centering
\begin{tabular}{ll}
Symbol & Meaning \\
\hline
$L$ & number of ligands in the fixed pool \\
$n_\text{families}$ & number of chemical families \\
$D$ & latent space dimension \\
$\sigma_\text{shape}$ & spread of a family's Gaussian \\
$d_\text{fam}$ & average distance between family centers \\
$\mu_S$ & mean number of ligands per sniff \\
$(\mu_{c,\ell},\sigma_{c,\ell})$ & concentration distribution of ligand $\ell$ \\
$\sigma_\text{noise}$ & observation noise on a ligand's coordinate \\
$\lambda$ & morphochemical length scale (Sec.~\ref{sec:kernel}) \\
\end{tabular}
\caption{Parameters defining the environment. Used in Sec.~\ref{sec:perfect}
to specify the regime in which the results are reported.}
\label{tab:env_params}
\end{table}

Given a sniff, the physics turns it into an activation pattern. A single
ligand activates receptor $r$ once its concentration exceeds the threshold,
$A_r=\Theta(\ln c-\ln c^*_{r,\ell})$ (Sec.~\ref{sec:mwc}). A sniff mixes
several ligands, which we combine into the receptor's drive,
\begin{equation}
\text{drive}_r = \sum_{\ell:\,M_\ell=1} \frac{c_\ell}{c^*_{r,\ell}},
\label{eq:drive}
\end{equation}
so that the receptor reaches half-activation at $\text{drive}_r=1$, which
recovers $c=c^*_{r,\ell}$ for a single ligand. The step $\Theta$ has zero
gradient almost everywhere, so for optimization we soften it into a
temperature-scaled sigmoid of the log-drive,
\begin{equation}
p(r\mid\text{sniff}) = \sigma\!\left(\frac{\ln\,\text{drive}_r}{T}\right),
\label{eq:soft_bin}
\end{equation}
which we evaluate through the numerically stable logsumexp
\begin{equation}
\ln\,\text{drive}_r = \ln\sum_{\ell:\,M_\ell=1} \exp\bigl(\ln c_\ell - \ln c^*_{r,\ell}\bigr),
\label{eq:logsumexp}
\end{equation}
since the ratios $c_\ell/c^*_{r,\ell}$ span many decades and a direct sum
can overflow or underflow. The temperature $T$ anneals from an initial
value to a small target. As $T\to0$, Eq.~\eqref{eq:soft_bin} recovers the
discrete activation $p(r\mid\text{sniff})\to A_r\in\{0,1\}$, so the array's
response is the binary pattern $A=(A_1,\dots,A_R)$.

Finally, we optimize the receptors. The environment is drawn once and held
fixed. The only free parameters are the receptor chemistry of
Sec.~\ref{sec:kernel}, $\mathbf v_u^{\pm}$, $E_\text{base}^{u,\pm}$,
$E_\text{max}^{u,\pm}$, which we update by backpropagation with the Adam
optimizer in PyTorch. Fig.~\ref{fig:env}\textbf{C} summarizes the loop:
each sniff yields a pattern $A$, repeated sniffs build the distribution
$p(A)$ and its entropy $H(A)$, and the receptor coordinates are moved to
maximize it.

\section{The perfect-array regime}
\label{sec:perfect}

A homomer-heteromer comparison is only meaningful once we know the
environment is not what limits the entropy: $H(A)$ can be capped by the
sniff's own statistics, by the receptor geometry's capacity to realize
distinct dichotomies, or by the optimizer's ability to reach and measure
$H(A)=R$ (each detailed below). We seek a regime where none of these three
mechanisms binds, the perfect-array regime, in which the homomer array
reaches its architectural ceiling, $H(A)=R$. Once none of the three limits
$H(A)$, the only thing left to set it is the receptor array's own
architecture, exactly what we want to isolate. This regime coincides with
a window in which $H(A)$ becomes insensitive to the precise value of each
environment parameter, which makes sense: once no mechanism binds, moving
a parameter within the window cannot change the outcome. Parameter
independence is therefore a consistency check on being in the
perfect-array regime, not a separate requirement, and we use it below as
one.

As we show in the main text, heteromer arrays evaluated in this same
window, with the same environment and the same fitting procedure, do not
reach $R$. Since the homomer array already does, throughout it, neither
the environment nor the optimization can explain a heteromer shortfall.
Any gap measured between heteromer and homomer arrays of the same $R$ is
therefore a property of the heteromeric array itself, its subunit-sharing
correlations (Sec.~\ref{sec:mwc}), not an artifact of the environment or of
how it is fit.

\begin{itemize}
\item \textbf{Statistical.} $A$ is a deterministic, noiseless function of
the sniff, so $H(A)\le H(\text{sniff})$, the entropy of the sniff itself,
which ligands are present and at what concentration: a deterministic map
cannot create entropy. If the sniff carries little entropy, too sparse or
too few ligands, $H(\text{sniff})$ binds before $R$ does.
\item \textbf{Geometric.} Reaching $H(A)=R$ needs $R$ receptors whose
regions of high affinity, ``balls'' of radius $\sim\lambda$ around each
receptor's position in latent space (Sec.~\ref{sec:kernel}), realize $R$
mutually decorrelated, half-active dichotomies of the ligand
pool~\cite{zwicker_receptor_2016}. Whether they can is a question of the
receptor family's expressive power in $D$ dimensions.
\item \textbf{Numerical.} The first two mechanisms are properties of the
physical model. The gradient-based procedure of Sec.~\ref{sec:env} can
separately fail for concrete numerical reasons: an infeasible batch size,
an intractable gradient, or a latent dimension $D$ too large to fit in
memory.
\end{itemize}

Table~\ref{tab:param_windows} lists, for each parameter, the range over
which $H(A)=R$ holds, and which mechanism bounds it.

\begin{table}[t]
\centering
\begin{tabular}{p{2.2cm}p{2.6cm}p{5.2cm}p{2.0cm}}
Parameter & Range & Justification & Mechanism \\
\hline
$\mu_S,\,L$ & not too small & Too sparse or too few ligands caps
$H(\text{sniff})$ below $R$. & Statistical \\[6pt]
$\rho=\dfrac{\sigma_\text{shape}\sqrt{D}}{\lambda}$ & $\sim[0.2,1]$ &
Below: ligands within a family share one receptor's ball. Above: each ball
covers at most one ligand. & Geometric \\[6pt]
$d_\text{fam}/\lambda$ & $\sim[0.5,1.5]$ & Below: neighbouring families'
balls overlap. Above: no ball bridges families. & Geometric \\[6pt]
$D$ & $[5,15]$ & We scale $\sigma_\text{shape}\propto1/\sqrt D$
(Sec.~\ref{sec:env}) to hold $\rho$ fixed, so $D$ is free of the geometric
window. Larger $D$ only makes the evaluation heavier (memory, parameter
count), so we take the smallest $D$ the geometric mechanism allows. &
Numerical \\
\end{tabular}
\caption{Parameter windows for the perfect-array regime
(Sec.~\ref{sec:perfect}).}
\label{tab:param_windows}
\end{table}

Classically, for linear classifiers, hyperplanes, any $D+1$ points in
general position in $\mathbb{R}^D$ can be split into any of their
$2^{D+1}$ dichotomies, but this guarantee is lost beyond $D+1$ points: this
maximal number of points a classifier can arbitrarily separate is its VC
dimension, $D+1$ for a hyperplane~\cite{cover_geometrical_1965}.

Our receptors are not hyperplanes: through the Gaussian kernel of
Sec.~\ref{sec:kernel}, each is a ``ball'' classifier responding to ligands
within $\lambda$ of its position, generally less constrained by $D$ than a
hyperplane since the induced feature space is effectively
infinite-dimensional. Increasing $D$ can therefore only make the
classifier at least as capable, never less.

\section{Bracketing the joint Shannon entropy}
\label{sec:entropy}

Exact evaluation of $H(A) = -\sum_A p(A)\log_2 p(A)$ requires the joint
probability of all $2^R$ activation patterns. Computing it from a batch of
$B$ sniffs needs, for each sniff, a length-$2^R$ tensor built by an outer
product across the $R$ per-receptor probabilities, so the computation
holds a $(B,2^R)$ tensor in memory regardless of how $B$ is chosen.
Reliable coverage further needs $B\sim2^R$ sniffs, one per pattern, at
which point the tensor has $(2^R)^2=2^{2R}$ entries, $2^{2R+2}$ bytes at
float32. At $R=15$ this is $2^{32}$ bytes, $4$ GiB, comfortable on an
80 GB A100. Each extra receptor quadruples it: by $R=18$ it is $2^{38}$
bytes, $256$ GiB, beyond any single GPU. We need an estimator that does not
build a $2^R$-sized tensor once $R$ grows past the high teens.

We first reformulate $H(A)$ as the entropy of the activation pattern
averaged over sniffs. Index the $B$ sniffs in a batch by $i=1,\dots,B$, and
let $C$ be the (uniform) label of which sniff produced a given sample of
$A$. $A_{i,r}\equiv p(r\mid\{c_\ell\})$ (Eq.~\eqref{eq:soft_bin}) is
receptor $r$'s activation probability for sniff $i$: conditional on sniff
$i$, the receptors are independent Bernoulli variables,
\begin{equation}
p(A\mid C=i) = \prod_{r=1}^R A_{i,r}^{\,A_r}\,(1-A_{i,r})^{\,1-A_r}.
\label{eq:pAcondC}
\end{equation}
Each parameter $A_{i,r}$ already combines every ligand of sniff $i$ through
the drive sum of Eq.~\eqref{eq:drive}, so component $i$ is the array's
response to that one sniff. Averaging over the batch,
\begin{equation}
p(A) = \frac1B\sum_i p(A\mid C=i),
\label{eq:pA_mixture}
\end{equation}
gives $p(A)$ as a mixture, in the statistical sense, of these $B$
components.

By definition of mutual information,
\begin{equation}
H(A) = H(A\mid C) + I(A;C).
\label{eq:mixture_decomp}
\end{equation}
The first term is cheap: it is the entropy of the product-Bernoulli
distribution of Eq.~\eqref{eq:pAcondC}, averaged over the batch,
\begin{equation}
H(A\mid C) \equiv H_\text{cond} = \frac{1}{B}\sum_i \sum_r h_2(A_{i,r}),
\qquad
h_2(a) = -a\log_2 a - (1-a)\log_2(1-a).
\label{eq:hcond}
\end{equation}
The second term, $I(A;C)$, is exactly as hard as $H(A)$ itself: it asks how
distinguishable the $B$ components are, which still needs the full joint
distribution. Kolchinsky and Tracey give certified upper and lower bounds
on it from pairwise affinities between components
alone~\cite{kolchinsky_estimating_2017}, avoiding the $2^R$ tensor above.
Writing $A_i\equiv(A_{i,1},\dots,A_{i,R})$ for sniff $i$'s vector of
receptor probabilities:

\begin{itemize}
\item \textbf{Lower bound.} The Bhattacharyya affinity between sniffs $i$
and $j$,
\begin{equation}
\log\text{BC}(i,j) = \sum_r \log\!\Bigl(\sqrt{A_iA_j}+\sqrt{(1-A_i)(1-A_j)}\Bigr),
\label{eq:BC}
\end{equation}
gives
\begin{equation}
H_\text{KT}^\text{low} = H_\text{cond} - \frac{1}{B}\sum_i \log_2\!\Bigl(\frac{1}{B}\sum_j \text{BC}(i,j)\Bigr).
\label{eq:KTlow}
\end{equation}
\item \textbf{Upper bound.} Replacing the Bhattacharyya affinity by minus
the Kullback-Leibler (KL) divergence between the same product-Bernoulli
components,
\begin{equation}
-\text{KL}(i\|j) = \sum_r \Bigl[A_i\log A_j + (1-A_i)\log(1-A_j)\Bigr] - \sum_r\Bigl[A_i\log A_i+(1-A_i)\log(1-A_i)\Bigr],
\label{eq:KL}
\end{equation}
gives
\begin{equation}
H_\text{KT}^\text{up} = H_\text{cond} - \frac{1}{B}\sum_i \log_2\!\Bigl(\frac{1}{B}\sum_j e^{-\text{KL}(i\|j)}\Bigr).
\label{eq:KTup}
\end{equation}
\end{itemize}

Both bracket terms have the form $H_\text{cond}+(\text{estimate of
}I(A;C))$. In general $I(A;C)\le H(C)$. Since $C$ is uniform over the $B$
sniffs in the batch, $H(C)=\log_2 B$, so $I(A;C)\le\log_2 B$.

$H_\text{KT}^\text{low} \le H(A) \le H_\text{KT}^\text{up}$ is a certified
bracket for any $B$. It tightens to zero width exactly as the array
approaches the perfect-array regime of Sec.~\ref{sec:perfect}: pushing
different sniffs to activate disjoint, sharply-defined sets of receptors,
through decorrelation and threshold sharpening, is the same objective the
optimization already pursues to maximize $H(A)$. The better the array
performs, the more precisely we can measure how well it performs.

\bibliographystyle{apsrev4-2}
\bibliography{smelling}

\end{document}
